\section{System Design}
This chapter provides design artifacts and data modeling notes.

\subsection{Unified Modeling Language (UML)}
\textbf{Planned diagrams for final submission:}
\begin{itemize}
  \item Activity diagram for local PDF workflow
  \item Use case diagram for user operations
  \item Class diagram for feature and service modules
  \item Sequence diagram for file import to export flow
\end{itemize}

\textit{Note: Diagram figures will be added in the next revision based on final feature lock.}

\subsection{Database Design}
The current system does not require an application database for core functionality because processing is local-first and session-based.

\subsubsection{Data Dictionary}
The core runtime entities are: input file, file metadata, operation configuration, and output artifact.

\subsubsection{Database Relationship Diagram}
A formal database relationship diagram is not applicable in the current architecture because no persistent relational database is used for primary PDF workflows.

\subsection{E-R Diagram}
An E-R diagram is optional for this project version and can be represented as a conceptual model of runtime entities rather than persisted tables.

\subsection{User Interface Design (System Layout)}
The interface follows a simple structure: tool selection, file selection, operation controls, preview or status area, and output download action. The design target is low learning effort for first-time users.
